% =================
% 致读者 — formal content notice (non-diegetic)
% Official page before the fictional immersion begins
% =================
\cleardoublepage
\thispagestyle{empty}
\begingroup
\setlength{\fboxsep}{8mm}
\begin{center}
{\color{sangue}\scshape\Large 致读者}\\[1.2em]
{\color{sangue}\rule{0.4\paperwidth}{0.6pt}}\\[2.2em]
\end{center}

\noindent 本书在其虚构部分中包含对以下内容的直接呈现：

\begin{itemize}
\itemsep0.4em
\item 机构性精神科暴力——包括电休克、脑叶切除术和精神病收容院监禁（巴巴塞纳收容院，1903—1994）；
\item 饥荒、死亡以及儿童和成人的集体埋葬（塞阿拉大旱，1915—1932）；
\item 详细的医学和尸检描述，使用直白的临床语言；
\item 死亡、哀伤、遗弃和长期隔离；
\item 第二人称叙事，直接召唤读者占据\textit{患者1,263}的位置，并带有模拟的阅读指示。
\end{itemize}

\vspace{1em}

\noindent \textbf{关于模拟指示:} 本书中的任何指示都不会产生真实的物理、化学或生物效应。“污染”是一种叙事手法。书页以普通纸张和油墨印刷，未涂敷任何物质。读者可在任何时候、在任何碎片处停止阅读，除故事本身暗示的后果外，不会有任何其他后果。

\vspace{1em}

\noindent \textbf{关于梦境、睡眠与症状:} 没有证据表明阅读本书会诱发特定的梦境、改变睡眠—觉醒周期，或产生任何生理症状。书中对身体反应的描写是一种形式手法，而非经验性断言。

\vspace{1em}

\noindent \textbf{关于临床代码与量表:} 虚构部分出现的诊断协议、ICD代码与评估量表均为文学建构。所提及的障碍是真实的临床状况，但将其应用于一个叙事性行动者并无任何医学依据。\textbf{本书不诊断读者}：它模拟诊断，作为一种虚构装置。

\vspace{1em}

\noindent \textbf{关于虚构与现实的区分:} 若阿金·安东尼奥·科雷亚、埃托尔·奥利维拉医生、卢西亚·门德斯医生以及怪物均为虚构。政府牲畜圈（1915—1932）和巴巴塞纳收容院（1903—1994）是有据可查的历史事实。完整区分见卷末《历史注记》（第 \pageref{nota-historica} 页）。

\vspace{1em}

\noindent \textbf{如果您对以上任一主题敏感:} 本书此刻可能不适合您。推迟阅读或放下本书并不可耻。书不应该是耐力测试。

\vspace{1em}

\noindent \textbf{如果您需要支持:} 生命价值关怀中心（CVV）188（巴西，24小时免费电话）或 \texttt{cvv.org.br}；988 自杀与危机生命线（美国）或 \texttt{988lifeline.org}；中国心理援助热线 12356；您当地的精神卫生服务；或您信任的人脉网络。

\vspace{2.2em}
\begin{center}
{\color{sangue}\rule{0.4\paperwidth}{0.6pt}}\\[1em]
{\small\itshape 档案员}
\end{center}
\endgroup
\null\vfill
\clearpage
